\documentclass[11pt,a4paper]{article}
\usepackage[utf8]{inputenc}
\usepackage[spanish]{babel}
\usepackage{geometry}
\geometry{left=2.5cm, right=2.5cm, top=3cm, bottom=3cm}
\usepackage{xcolor}
\usepackage{graphicx}
\usepackage{hyperref}
\usepackage{booktabs}
\usepackage{tabularx}
\usepackage{titlesec}
\usepackage{enumitem}
\usepackage{amssymb}

% Definición de colores para las propuestas de diseño
\definecolor{greennewen}{HTML}{1B5E20}
\definecolor{goldnewen}{HTML}{F9A825}
\definecolor{graynewen}{HTML}{333333}
\definecolor{tradgreen}{HTML}{2E4A1E}
\definecolor{cremabg}{HTML}{F5ECD7}
\definecolor{claybrown}{HTML}{8D5524}

\titleformat{\section}{\Large\bfseries\color{greennewen}}{}{0em}{}[\titlerule]
\titleformat{\subsection}{\large\bfseries\color{graynewen}}{}{0em}{}

\begin{document}

\begin{center}
    {\Huge \bfseries Plan de Desarrollo}\\[0.5cm]
    {\LARGE \textcolor{greennewen}{Plataforma Web Comercial Digital Transaccional}}\\[0.2cm]
    {\Large Cliente: Cooperativa Mapuche Newen Leufu}\\[0.8cm]
    
    \textbf{Autor y Consultor de Ingeniería:}\\
    \textbf{José González Cortés} -- Ingeniero Civil Industrial\\
    Email: \href{mailto:mixto@mixto.cl}{mixto@mixto.cl} | Tel: +56 958751274 | Concepción, Chile\\[0.5cm]
    \rule{\textwidth}{0.5pt}
\end{center}

\vspace{0.3cm}

\section{1. Propósito y Contexto del Proyecto}
La Cooperativa de Turismo y Desarrollo Territorial Mapuche Newen Leufu requiere transformar su sitio web actual (informativo y estático) en un \textbf{canal activo de ventas y reservas digitales}. El objetivo principal es eliminar la fricción de la reserva manual vía WhatsApp o correo electrónico, ofreciendo cotización automática y procesamiento de pagos online seguros para turistas nacionales e internacionales, destacando sus 4 locaciones operativas: Trawünko, Kiñewün, Secretos de la Ñaña y Kura Ruka.

\section{2. Arquitectura de la Plataforma (Módulos Funcionales)}
La plataforma web de Newen Leufu se organiza en los siguientes módulos funcionales obligatorios:

\begin{itemize}[label=\textcolor{goldnewen}{$\bullet$}]
    \item \textbf{Front Page (Página de Inicio):} Hero Section de alto impacto con propuesta de valor bilingüe. Sección "¿Quiénes Somos?" para transmitir la misión y el concepto de turismo regenerativo mapuche. Preview de las 4--5 experiencias principales. Sección de locaciones con mapa/galería, testimonios reales de clientes, formulario de contacto, enlace directo a WhatsApp Business y mapa embebido.
    \item \textbf{Fichas Individuales de Servicio:} Páginas dedicadas bilingües con video responsivo embebido, carrusel de 5--8 fotos, narrativa de la experiencia, tabla técnica completa (duración, precio, mínimo/máximo de participantes, dificultad, idioma guía), botón de reserva directa, reseñas específicas e integración con TripAdvisor.
    \item \textbf{Catálogo de Servicios:} Vista en grilla con filtros dinámicos (categorías, duración, precio, locación) y ordenación (precio, popularidad).
    \item \textbf{Cotizador Personalizado (4 pasos):} Selección de experiencias, número de participantes (adultos, estudiantes y niños), calendario de disponibilidad en tiempo real, servicios adicionales (fotografía profesional) y resumen de precios en CLP y USD.
    \item \textbf{Carrito y Checkout Transaccional:} Resumen de la reserva con captura de datos del turista, política de cancelación visible y selección de pasarelas de pago.
    \item \textbf{Panel de Administración de Reservas:} Panel interno simplificado para que la Cooperativa gestione y confirme reservas sin conocimientos técnicos avanzados.
\end{itemize}

\section{3. Estrategia Técnica y Solución de Rendimiento}
Para garantizar la exigencia de \textbf{tiempos de carga inferiores a 3 segundos en móviles 4G} y evitar la sobrecarga ("bloat") de un WordPress tradicional con múltiples plugins de reserva y traducción, se propone una \textbf{arquitectura desacoplada (Headless)} de alto rendimiento:

\begin{itemize}[label=\textcolor{greennewen}{$\bullet$}]
    \item \textbf{Frontend (SPA):} Desarrollado en Vue 3 con TypeScript y Tailwind CSS v4. Garantiza interactividad fluida, cambio de idioma instantáneo (Vue I18n) y transiciones sin recarga de navegador.
    \item \textbf{Backend \& CMS:} Headless CMS / API REST segura en Node.js que permite gestionar el catálogo de experiencias, la disponibilidad de fechas y recibir las reservas.
    \item \textbf{Pasarelas de Pago:} Integración con Webpay Plus de Transbank para pagos nacionales en CLP (Débito/Crédito) y Stripe / PayPal para turistas internacionales en USD/EUR.
    \item \textbf{Seguridad y Despliegue:} Certificado SSL (HTTPS obligatorio), estándares PCI-DSS de seguridad en transacciones y correos transaccionales seguros por SMTP (SendGrid o Mailchimp).
\end{itemize}

\subsection{3.1. Justificación Técnica: Análisis del Sitio Actual vs. Nueva Plataforma}
Para fundamentar la necesidad del cambio tecnológico, se realizó un análisis técnico y de rendimiento en tiempo real sobre el sitio web actual (\url{https://newenleufu.cl}), identificando deficiencias de infraestructura críticas que merman las ventas:

\begin{itemize}[itemsep=2pt]
    \item \textbf{Monolito WordPress Ineficiente:} El sitio actual corre sobre WordPress con consultas constantes a la base de datos PHP/MySQL. Esto genera un \textbf{TTFB (Time to First Byte) de 956 ms} (casi 1 segundo de retraso en el servidor antes de enviar datos), lo cual penaliza severamente el posicionamiento SEO en Google (Core Web Vitals).
    \item \textbf{Carga Lenta del Documento:} Descargar únicamente el archivo HTML inicial toma \textbf{1.39 segundos}, una latencia sumamente alta considerando conexiones móviles 4G inestables en zonas rurales o turistas en el extranjero con mayor distancia de ruteo.
    \item \textbf{Ventaja Radical de la Nueva Plataforma (Vue 3 + Cloudflare Pages):}
    \begin{itemize}
        \item \textbf{TTFB menor a 25 ms} (una mejora de más del \textbf{3.800\%} en la respuesta inicial) al entregar archivos compilados estáticos directamente desde la red global de borde (Edge CDN) de Cloudflare.
        \item \textbf{Carga instantánea de la UI en menos de 150 ms}, lo que ofrece una sensación de velocidad nativa, incrementando significativamente las conversiones de reserva.
    \end{itemize}
\end{itemize}

\section{4. Enfoque de Diseño: Unificación de Propuestas Estéticas}
La plataforma comercial digital de Newen Leufu integra de manera única y reactiva las dos propuestas estéticas, permitiendo al usuario alternar entre ellas en tiempo real mediante un conmutador de temas inteligente, sin necesidad de mantener dos bases de código separadas:

\begin{enumerate}
    \item \textbf{Propuesta 1: "Raíces Vivas" (Enfoque Tradicional e Identidad Territorial):}
    Optimizado para transmitir la espiritualidad, herencia e identidad del pueblo mapuche y el Wallmapu.
    \begin{itemize}[itemsep=2pt]
        \item \textbf{Paleta de colores:} Verde musgo oscuro (\textcolor{tradgreen}{\rule{1.3ex}{1.3ex}}\,\#2E4A1E o \textcolor{greennewen}{\rule{1.3ex}{1.3ex}}\,\#1B5E20), crema de fondo suave (\textcolor{cremabg}{\rule{1.3ex}{1.3ex}}\,\#F5ECD7) y acentos de tierra arcilla (\textcolor{claybrown}{\rule{1.3ex}{1.3ex}}\,\#8D5524).
        \item \textbf{Tipografía:} Serif de estilo editorial y artesanal (Cormorant Garamond) para títulos principales, combinada con Lato para textos corporativos.
        \item \textbf{Elementos visuales:} Texturas radiales de grano de madera, separadores winkul geométricos tradicionales mapuches, y secciones dedicadas de sabiduría ancestral y cosmovisión.
    \end{itemize}
    
    \item \textbf{Propuesta 2: "Experiencia Moderna" (Enfoque Comercial y Conversión Rápida):}
    Optimizado para la máxima claridad, simplicidad de uso en dispositivos móviles y agilidad transaccional.
    \begin{itemize}[itemsep=2pt]
        \item \textbf{Paleta de colores:} Verde institucional vibrante (\textcolor{greennewen}{\rule{1.3ex}{1.3ex}}\,\#1B5E20), fondos blancos limpios y acentos de oro moderno (\textcolor{goldnewen}{\rule{1.3ex}{1.3ex}}\,\#F9A825).
        \item \textbf{Tipografía:} Sans-serif de alta legibilidad y estilo moderno (Poppins).
        \item \textbf{Elementos visuales:} Tarjetas de bordes redondeados modernos, espaciados amplios y minimalistas, y composición limpia de imágenes de alta resolución sin ornamentaciones tradicionales.
    \end{itemize}
\end{enumerate}

\newpage

\section{5. Plan de Desarrollo Detallado y Fases de Ingeniería}

Para garantizar un desarrollo ágil y la correcta implementación del alcance, el proyecto se estructurará en tres fases de desarrollo progresivas, con entregas funcionales constantes:

\begin{table}[h!]
    \centering
    \renewcommand{\arraystretch}{1.3}
    \begin{tabularx}{\textwidth}{|l|X|}
    \hline
    \textbf{Fase de Desarrollo} & \textbf{Hitos y Entregables} \\
    \hline
    \textbf{Fase 1: Frontend (24 horas)} & 
    -- \textbf{Setup y Estructura:} Configuración de Vue 3 + TS + Tailwind CSS v4, enrutamiento responsivo y layouts de la aplicación.\\
    & -- \textbf{Vistas Públicas:} Implementación de la Página de Inicio, sección bilingüe de "Quiénes Somos", carrusel TripAdvisor y formularios de contacto.\\
    & -- \textbf{Catálogo de Experiencias:} Desarrollo de la grilla interactiva de 6 experiencias bilingües oficiales, con barra de filtros por categorías y locaciones. \\
    \hline
    \textbf{Fase 2: Lógica (16 horas)} & 
    -- \textbf{Cotizador e Internacionalización:} Lógica del cotizador interactivo de 4 pasos (seleccionador de fecha con calendario de disponibilidad e inline pricing en CLP/USD) y soporte para traducción nativa (Español / Inglés).\\
    & -- \textbf{Checkout y Pasarelas de Pago:} Desarrollo del carrito, políticas de cancelación, modal integrado de checkout y simulación interactiva de pago. \\
    \hline
    \textbf{Fase 3: Integración (8 horas)} & 
    -- \textbf{Backend y CMS:} Creación de la base de datos de reservas, API REST y el panel administrativo simplificado (CMS) para bloquear fechas y gestionar reservas.\\
    & -- \textbf{Despliegue y SEO:} Configuración de correos transaccionales por SMTP, metatags de SEO, Sitemap, analíticas de conversión y despliegue automatizado continuo en Cloudflare Pages. \\
    \hline
    \end{tabularx}
\end{table}

\section{6. Cotización de Horas y Presupuesto}

La estimación presupuestaria se realiza bajo un enfoque simplificado de desarrollo por componentes, asignando un estándar de \textbf{8 horas de ingeniería especializada por cada módulo funcional principal} de la plataforma, cubriendo el 100\% del alcance técnico, integraciones de pasarelas, base de datos de reservas y despliegue descritos en este documento:

\vspace{0.3cm}
\begin{center}
    \renewcommand{\arraystretch}{1.3}
    \begin{tabularx}{0.95\textwidth}{|X|r|}
    \hline
    \textbf{Módulo Funcional / Componente a Desarrollar} & \textbf{Horas Estimadas} \\
    \hline
    Módulo 1: Front Page (Inicio, Quiénes Somos, Testimonios, Contacto, WhatsApp) & 8 hrs \\
    Módulo 2: Fichas de Experiencia Individuales (Galería, Video, Datos Técnicos) & 8 hrs \\
    Módulo 3: Catálogo de Servicios con Filtros Reactivos y Ordenación & 8 hrs \\
    Módulo 4: Cotizador Personalizado e Interactivo de 4 pasos & 8 hrs \\
    Módulo 5: Carrito de Compras y Checkout Transaccional Unificado & 8 hrs \\
    Módulo 6: Panel de Administración y Back-End de Gestión (Disponibilidad y CMS) & 8 hrs \\
    \hline
    \textbf{Total General de Horas Estimadas para la Plataforma Comercial} & \textbf{48 hrs} \\
    \hline
    \end{tabularx}
\end{center}

\vspace{0.3cm}
\noindent \textit{* Esta cotización de 48 horas de ingeniería de software cubre de extremo a extremo la implementación del alcance completo, garantizando una arquitectura desacoplada bilingüe de alto rendimiento, libre de las ineficiencias de un CMS monolítico tradicional como WordPress y lista para desplegarse globalmente con entrega continua en Cloudflare Pages.}

\end{document}